% ---- Abstract in Khmer ----
% NOTE: This is a DRAFT machine-assisted Khmer translation of the English abstract.
%       Please PROOFREAD and adjust wording/terminology before submission.
%       English/Latin terms are wrapped in \en{...} so they render in the Latin font
%       (Noto Sans Khmer has no Latin glyphs).
\chapter*{\km{សេចក្ដីសង្ខេប}}
\addcontentsline{toc}{chapter}{Abstract (Khmer)}

\begin{khmer}
ការដាក់ពិន្ទុចម្លើយខ្លីដោយស្វ័យប្រវត្តិ (\en{ASAG}) បានរីកចម្រើនយ៉ាងឆាប់រហ័សសម្រាប់ភាសាដែលមានធនធានច្រើន
ប៉ុន្តែភាសាខ្មែរ ដែលជាភាសាផ្លូវការរបស់ប្រទេសកម្ពុជា និងនិយាយដោយប្រជាជនជាង ១៧ លាននាក់
មិនទាន់មានមូលដ្ឋានគោល (\en{benchmark}) សម្រាប់ \en{ASAG} ដែលបានបោះពុម្ពផ្សាយពីមុនមកនោះទេ។
និក្ខេបបទនេះបង្កើតមូលដ្ឋានគោល \en{ASAG} ភាសាខ្មែរដំបូងគេ ដែលអាចផលិតឡើងវិញបាន៖
សារពើភ័ណ្ឌនៃចម្លើយខ្លីចំនួន ១,១៨៤ ដែលដាក់ពិន្ទុដោយគ្រូ លើមុខវិជ្ជាចំនួន ៤
ព្រមទាំងបណ្ដុំទិន្នន័យសម្អាតចំនួន ៣ និងខួរផលិត (\en{pipeline}) រួមមួយ
ដែលប្រៀបធៀបគ្រួសារគំរូចំនួន ៤៖ គំរូបុរាណ (លក្ខណៈ \en{TF-IDF} ជាមួយ \en{Support-Vector Regression}),
បណ្ដាញ \en{BiLSTM} ថ្នាក់តួអក្សរ ជាមួយ \en{Attention}, គំរូ \en{Transformer} ពហុភាសា (\en{mBERT, XLM-R, GTE}),
និងគំរូភាសាធំ (\en{LLM}) ដែលបាន \en{fine-tune} ដោយ \en{QLoRA} (\en{Qwen})។

ដោយសារចំណងជើងផ្ដោតលើ ការពន្យល់បាន (\en{Explainability}) និក្ខេបបទនេះក៏ធ្វើការសិក្សាការពន្យល់
ឆ្លងគ្រួសារគំរូ ដែលវាស់វែង ភាពស្មោះត្រង់ (\en{faithfulness}) នៃការពន្យល់ (\en{ERASER comprehensiveness},
\en{sufficiency} និង \en{AOPC}) និងភាពសមហេតុផល (\en{plausibility}) ធៀបនឹងចម្លើយយោង
មិនមែនគ្រាន់តែបង្ហាញការបន្លិចពាក្យសំខាន់ៗប៉ុណ្ណោះទេ។

លទ្ធផលសំខាន់ៗមានបី។ ទីមួយ គ្រួសារគំរូទាំងបួនមានលទ្ធផល ប្រហាក់ប្រហែលគ្នា
លើ \en{Quadratic Weighted Kappa (QWK)} គឺ ០.៧៩៥, ០.៨៤៥, ០.៨២០ និង ០.៨៤៣ (ស្ថិតក្នុងចន្លោះតូចចង្អៀត ០.០៥)
រីឯ \en{LLM} ដែលបាន \en{fine-tune} គឺល្អបំផុតយ៉ាងច្បាស់លាស់សម្រាប់រង្វាស់ការដាក់ឱ្យប្រើ
(ត្រូវគ្នាជាក់លាក់ ៦៦\% និងត្រូវគ្នាក្នុងរង្វង់មួយពិន្ទុ ៨៣\%)។
ទីពីរ ការពន្យល់បែបការផ្ដល់សារៈសំខាន់ដល់ពាក្យតាមវិធី \en{Leave-One-Out (LOO)}
មានភាពស្មោះត្រង់ ជឿជាក់បាន សម្រាប់គ្រប់គ្រួសារគំរូ ដែលបានផ្ទៀងផ្ទាត់ដោយ
\en{ERASER comprehensiveness} និង \en{sufficiency} ធ្វើឱ្យ \en{LOO} ជាវិធីពន្យល់រួមមួយ
ឯករាជ្យពីគំរូ (\en{model-agnostic}) សម្រាប់សសរស្ដម្ភទាំងបួន។
ទីបី ក្រោមការវាយតម្លៃបែបទុកសំណួរចេញ (\en{question-held-out}) \en{QWK} របស់គំរូបុរាណធ្លាក់ចុះពី ០.៧៦ ទៅ ០.៣៥
ដែលបង្ហាញថាពិន្ទុ \en{ASAG} លើបណ្ដុំសំណួរតូច ត្រូវបានបំប៉ោងដោយ ការលេចធ្លាយសំណួរ (\en{question leakage})។
ការរួមចំណែកផ្នែកការពន្យល់ត្រូវបានសម្រេចជា គំរូបណ្ដាញតាមអ៊ីនធឺណិត (\en{web prototype}) សម្រាប់គ្រូ
ដែលផ្ដល់ពិន្ទុ ការពន្យល់ពាក្យសំខាន់ និងមតិកែលម្អ។ កូដ ទិន្នន័យ និងលទ្ធផលទាំងអស់ត្រូវបានចេញផ្សាយ
ដើម្បីភាពអាចផលិតឡើងវិញ។
\end{khmer}

\vspace{0.3cm}
\noindent\textbf{\km{ពាក្យគន្លឹះ}:} \km{ការដាក់ពិន្ទុចម្លើយខ្លីដោយស្វ័យប្រវត្តិ; បញ្ញាសិប្បនិមិត្តពន្យល់បាន;
ភាពស្មោះត្រង់; ដំណើរការភាសាខ្មែរ; ភាសាធនធានតិច;} \en{Transformers; Quadratic Weighted Kappa.}
\clearpage
